\section{Related Work}

\paragraph{Ambiguity in language and question answering.} Ambiguity is a design feature of efficient communication rather than a defect \citep{piantadosi2012communicative}, and listeners resolve most of it incrementally from context \citep{ferreira2008ambiguity, sedivy1999achieving}. When context does not suffice, interlocutors ground the intended reading through dialogue \citep{clark1991grounding, clark1996using}. Text-only benchmarks operationalize the phenomenon as questions with several plausible answers: AmbigQA \citep{min2020ambigqa} asks models to surface all of them, and follow-up work studies when models should ask instead of answering \citep{stengeleskin2023chicken}. In these resources the ambiguity lives in world knowledge or in retrieval. \bench{} moves it into the visual referent of a fixed surface question.

\paragraph{Referential ambiguity in images.} RAcQUEt \citep{testoni2025racquet} shows that image-grounded models answer referentially ambiguous questions about one referent without acknowledging the others, a behavior the authors link to overconfidence and to social bias. ClearVQA \citep{jian2025clearvqa} trains models to ask clarifying questions on ambiguous image--question pairs, MUCAR \citep{wang2025mucar} extends the setting across languages, and earlier work probed uncertainty acknowledgment more broadly \citep{liu2023afraid, pezzelle2023dealing}. In all of these, the competing referents are co-present in a single frame, so silent commitment can persist even when every candidate is perceived. In video, the readings are distributed over time and must be computed through event individuation and identity tracking \citep{zacks2007event}, so the same surface behavior can hide a different fault. Our diagnostics separate the two, and \S\ref{sec:diagnostics} locates the video fault in computing the readings.

\paragraph{Video benchmarks and hallucination.} Standard video question-answering benchmarks such as MVBench \citep{li2024mvbench} and Video-MME \citep{fu2024videomme} assign one correct answer per question, so a model that commits to a single reading is never penalized. Video hallucination benchmarks measure fabricated content \citep{liu2025vidhalluc, choong2025vidhal}, following image-side analyses of object hallucination \citep{rohrbach2018object}; question ambiguity itself has been treated as a nuisance variable to remove \citep{bhattacharya2019why}. The failure \bench{} targets is complementary to both: nothing the model says is fabricated, and the harm lies in the $K-1$ readings it never considered.

Table~\ref{tab:comparison} summarizes the contrast with the closest resources along the axes that matter for hallucinated commitment: modality, temporal distribution of the alternatives, exhaustive gold reading lists, and an interactive protocol that lets asking pay off.
